% Source: https://artofproblemsolving.com/community/c3h1669858_a_little_help_with_this_please
% Stable identity: aops-topic-1669858-q01
\textbf{Domain.} We need $x+3\ge0$, $-x-1\ge0$ and $-x^2-4x-3=-(x+1)(x+3)\ge0$, i.e.\ $-3\le x\le-1$.

\textbf{Substitution.} Put $u=\sqrt{x+3}\ge0$ and $v=\sqrt{-x-1}\ge0$. Then $u^2+v^2=2$ and $uv=\sqrt{(x+3)(-x-1)}=\sqrt{-x^2-4x-3}$. With $w=u-v$ we have $w^2=2-2uv$, so $uv=1-\dfrac{w^2}{2}$, and the inequality becomes
\[
w<2-\dfrac{w^2}{2}\iff w^2+2w-4<0\iff -1-\sqrt5<w<-1+\sqrt5 .
\]
Since $w\ge-v\ge-\sqrt2>-1-\sqrt5$, the condition is just $u-v<\sqrt5-1$.

\textbf{Solving $u-v<\sqrt5-1$.} On $[-3,-1]$, $u$ increases and $v$ decreases with $x$, so $u-v$ is strictly increasing, from $-\sqrt2$ at $x=-3$ to $\sqrt2$ at $x=-1$; as $\sqrt5-1<\sqrt2$, there is exactly one $x_0$ with $u-v=\sqrt5-1$, and the solution set is $-3\le x<x_0$. At $x_0$: $w^2=6-2\sqrt5$, $u+v=\sqrt{2+2uv}=\sqrt{4-w^2}=\sqrt{2\sqrt5-2}$, and
\[
2x_0+4=u^2-v^2=(u-v)(u+v)=\sqrt{(6-2\sqrt5)(2\sqrt5-2)}=\sqrt{16\sqrt5-32}=4\sqrt{\sqrt5-2}.
\]
So $x_0=2\sqrt{\sqrt5-2}-2\approx-1.028$, and the solution is
\[
\boxed{-3\le x<2\sqrt{\sqrt5-2}-2}.
\]
